\chapter{Introduzione ai Big Data}

\section{Che cosa sono i Big Data?}
Il termine \textbf{Big Data} è generalmente associato a enormi quantità di dati ("petabytes") raccolte in grandi collezioni. Tuttavia, il concetto non si limita alla sola quantità, ma coinvolge diverse dimensioni fondamentali:
\begin{itemize}
    \item \textbf{Storage}: i dati devono essere memorizzati utilizzando approcci efficienti ed efficaci. Spesso non si utilizzano i tradizionali database relazionali, ma ci si affida a \textit{distributed file system} o database NoSQL.
    \item \textbf{Analytics}: l'obiettivo finale è estrarre conoscenza e valore dai dati attraverso metodologie di analisi.
    \item \textbf{Technologies}: sono necessarie nuove tecnologie e framework per elaborare e analizzare questi dati in modo scalabile.
\end{itemize}

Inizialmente, i framework per i Big Data sono stati sviluppati dalle grandi aziende di internet (le "Big Tech" come Google e Meta). Oggi, tuttavia, la necessità di analizzare enormi moli di dati riguarda qualsiasi settore, dalle banche agli impianti industriali.

\begin{remark}
\textbf{Sfruttare i dati esistenti per nuovi task}\\
Il professore sottolinea che un buon \textit{data scientist} o \textit{data engineer} deve saper utilizzare i dati già a disposizione per risolvere nuovi problemi. Ad esempio, Google ha iniziato a raccogliere le query degli utenti per fornire risposte migliori, ma ha poi utilizzato gli stessi dati per fare \textit{nowcasting} (previsione a brevissimo termine), stimando in tempo reale quante persone fossero affette da influenza in una determinata area.
\end{remark}

\section{Fonti e Generatori di Big Data}
Esistono diverse fonti in grado di generare enormi quantità di dati, che tipicamente vengono poi convogliate e memorizzate in un singolo Data Center. Le principali categorie includono:

\begin{enumerate}
    \item \textbf{User Generated Content (Web \& Mobile)}: contenuti generati dagli utenti su piattaforme come Facebook, Instagram, Yelp, TripAdvisor, Twitter e YouTube.
    
% FIGURA DA INSERIRE QUI
% Fonte: Slide 6 (BD - Cap 1)
% Descrizione: Dispositivi mobile, loghi dei social media e un carretto pieno di icone social.

    \item \textbf{Health and scientific computing}: dati ad alta risoluzione provenienti dal settore medico, immagini satellitari (utili ad esempio per identificare aree allagate o bruciate), microscopi e telescopi.
    
% FIGURA DA INSERIRE QUI
% Fonte: Slide 6 (BD - Cap 1)
% Descrizione: Immagini di un microscopio, un telescopio e un satellite spaziale.


    \item \textbf{Log files}: file di testo generati da web server o sistemi automatici che tracciano eventi e accessi.
    
    \item \textbf{Internet Of Things (IoT)}: reti di sensori (\textit{Sensor networks}), RFID e \textit{smart meters}. Anche se ogni singolo sensore invia una quantità minuscola di dati, l'aggregazione di miliardi di sensori genera moli di dati enormi.
\end{enumerate}

\section{Le "V" dei Big Data}
Il contesto dei Big Data è tradizionalmente descritto attraverso le \textbf{3V}:
\begin{itemize}
    \item \textbf{Volume (Scale of data)}: la scala dimensionale dei dati. L'unità di misura tipica entra nell'ordine dei Petabyte (PB). Intorno al 2012, Google processava oltre 100 PB di dati al giorno, mentre Facebook oltre 10 PB.
    \item \textbf{Variety (Different forms of data)}: l'eterogeneità dei dati in ingresso. Un sistema Big Data dovrebbe poter gestire testo, immagini, grafi e log.
    \item \textbf{Velocity (Analysis of streaming data)}: la velocità con cui i dati vengono generati e la necessità di analizzare flussi di dati continui (stream) elaborando milioni di eventi in pochi millisecondi.
\end{itemize}

A queste, nel tempo, sono state aggiunte altre due dimensioni:
\begin{itemize}
    \item \textbf{Veracity (Uncertainty of data)}: la qualità dei dati.
    \item \textbf{Value}: la capacità di estrarre informazioni utili dal dato.
\end{itemize}

\begin{remark}
\textbf{Veracity e qualità del dato}\\
Il professore ricorda il principio \textit{"Garbage in, garbage out"}: se l'input è rumoroso o di scarsa qualità, anche l'output degli algoritmi di Machine Learning sarà pessimo. Inoltre, i dati grezzi raramente sono pronti per l'analisi: in una tipica pipeline di Data Science, dal 60\% al 70\% del tempo è dedicato alla pulizia del dato (\textit{data cleaning}) e alla sua trasformazione in vettori o matrici.
\end{remark}

\section{Framework Big Data: Motivazioni e Sfide}

\subsection{Perché abbiamo bisogno di nuovi framework?}
Nel passato avevamo già a disposizione database relazionali efficienti, sistemi NoSQL e l'High-Performance Computing (HPC). Tuttavia, questi sistemi mostrano dei limiti con moli di dati eccezionali. Ad esempio, per indicizzare 10 miliardi di pagine web (da circa 20 KB l'una), un \textit{full scan} sequenziale su un disco rigido tradizionale o l'estrazione a blocchi nella memoria RAM (come nel classico \textit{Data Mining}) richiederebbe un tempo inaccettabile. 

Occorreva passare da un'architettura a nodo singolo (\textit{single-node architecture}) a un approccio fortemente distribuito.

\subsection{Il principio della Data Locality}
Il concetto base sviluppato da Google per i framework Big Data è la \textbf{Data Locality}. 
\begin{definition}[Data Locality]
Invece di inviare enormi quantità di dati attraverso la rete verso i server che contengono il codice da eseguire, si invia il codice (che è di dimensioni molto ridotte) direttamente ai server in cui i dati sono fisicamente memorizzati.
\end{definition}
Inviare troppi dati sulla rete crea dei colli di bottiglia (\textit{network bottleneck}) che rendono l'applicazione lentissima. 

\subsection{Architettura dei Data Center e Scalabilità}
I sistemi Big Data utilizzano \textbf{commodity hardware}, ovvero server standard, economici e facilmente sostituibili in caso di guasto. Questi server sono raggruppati in rack e interconnessi tramite switch di rete. 

\begin{remark}
\textbf{High-Performance Computing (HPC) vs Big Data Framework}\\
L'HPC è indicato per compiti intensivi dal punto di vista computazionale (es. calcolare miliardi di operazioni in virgola mobile al secondo su un set di dati relativamente piccolo) e si basa sulla \textit{scalabilità verticale}. I framework Big Data (come Hadoop o Spark) sono invece adatti a carichi di lavoro \textbf{data-intensive} e si basano sulla \textbf{scalabilità orizzontale}, dove la metrica principale è quanti dati il sistema riesce ad analizzare in un certo lasso di tempo.
\end{remark}

\subsection{Le sfide tecniche}
Secondo le slide, le principali sfide affrontate dai framework Big Data sono:
\begin{itemize}
    \item Calcolo distribuito e parallelizzazione;
    \item Storage distribuito di dataset di dimensioni enormi (Terabyte, Petabyte);
    \item Gestione dei guasti dei nodi (\textit{Node failure management});
    \item Colli di bottiglia della rete (\textit{Network bottleneck});
    \item Formati di input eterogenei (\textit{Data diversity \& heterogeneity}).
\end{itemize}
Il vantaggio dell'utilizzo di framework come Spark è che lo sviluppatore non deve occuparsi direttamente dello scheduling dell'applicazione, della gestione della rete o dei fallimenti hardware, poiché vengono gestiti automaticamente dal framework.

\subsection{Il problema tipico nei Big Data}
Non tutti i problemi possono essere risolti efficientemente con un approccio Big Data (ad esempio, quando è necessaria una visione "globale" immediata dei dati per prendere una decisione). 
Inoltre, le funzioni applicate ai dati distribuiti devono soddisfare specifici vincoli matematici: devono essere \textbf{commutative} e \textbf{associative}.

Un problema tipico che ben si adatta a questi framework segue questo schema operativo:
\begin{enumerate}
    \item Iterare su una quantità enorme di record/oggetti;
    \item Estrarre l'informazione di interesse da ogni record (\textit{Map});
    \item Aggregare i risultati intermedi (\textit{Reduce});
    \item Generare l'output finale / risultato globale.
\end{enumerate}